% =====================================================================
\section{Giới thiệu}
% =====================================================================

% ---------------------------------------------------------------------
\begin{frame}{Encoder-only: nền tảng của các bài toán hiểu ngôn ngữ}
\vspace{0.2em}
\begin{columns}[T]
\begin{column}{0.54\textwidth}
  \textbf{Encoder-only là gì?}
  \begin{itemize}\itemsep3pt
    \item Đọc câu theo \emph{cả hai chiều}, sinh ra biểu diễn ngữ nghĩa
          cho từng từ -- không sinh văn bản.
    \item Là lựa chọn mặc định cho phân loại văn bản, gán nhãn, hỏi đáp
          trích xuất, tìm kiếm ngữ nghĩa.
    \item BERT và RoBERTa vẫn đang được dùng rộng rãi trong thực tế.
  \end{itemize}
\end{column}

\begin{column}{0.43\textwidth}
  \begin{block}{Vấn đề}
    Trong khi các mô hình sinh văn bản (LLM) tiến hóa liên tục, các mô
    hình dòng Encoder-only đã \emph{chững lại nhiều năm}, chỉ được cập
    nhật tiếp từ cuối năm 2024.
  \end{block}
\end{column}
\end{columns}

\vspace{0.5em}
\centering
% Mốc thời gian hai dòng mô hình. Trục hoành: 1 đơn vị = 1 năm, gốc = 01/2018,
% vị trí mốc tính theo tháng phát hành thực tế. Nhãn xếp so le hai tầng để các
% mốc gần nhau không đè lên nhau. Dòng Encoder dùng nét đứt cho quãng 2021-2024.
\fitpic{\begin{tikzpicture}[font=\tiny,x=1.62cm,y=1.32cm]

  % ---- dòng mô hình sinh văn bản: nhiều nhóm nghiên cứu, năm nào cũng có mốc
  \node[anchor=east,text=vnSteel,font=\scriptsize\bfseries] at (-0.3,0.72) {LLM};
  \draw[vnSteel,thick,-{Latex[length=2mm]}] (-0.2,0.72) -- (7.55,0.72);
  % tầng gần trục
  \foreach \x/\lb in {0.42/GPT-1, 2.33/GPT-3, 4.25/PaLM, 5.20/GPT-4}{
    \fill[vnSteel] (\x,0.72) circle (1.4pt);
    \node[above=1pt,text=vnSteel] at (\x,0.72) {\lb};
  }
  % tầng xa trục
  \foreach \x/\lb in {1.08/GPT-2, 3.94/Gopher, 4.87/ChatGPT, 6.25/{Llama 3},
                      7.05/{DeepSeek-R1}}{
    \fill[vnSteel] (\x,0.72) circle (1.4pt);
    \draw[vnSteel!45,line width=0.3pt] (\x,0.72) -- (\x,0.94);
    \node[above=0pt,text=vnSteel] at (\x,0.94) {\lb};
  }

  % ---- trục năm ở giữa
  \foreach \x/\yr in {0/2018,1/2019,2/2020,3/2021,4/2022,5/2023,6/2024,7/2025}{
    \node[text=vnInk!65] at (\x,0.30) {\yr};
  }

  % ---- dòng Encoder: liền nét -> đứt quãng -> liền nét trở lại
  \node[anchor=east,text=vnNavy,font=\scriptsize\bfseries] at (-0.3,-0.02) {Encoder};
  \draw[vnNavy,thick]        (-0.2,-0.02) -- (3.88,-0.02);
  \draw[vnNavy,thick,dashed] (3.88,-0.02) -- (6.80,-0.02);
  \draw[vnNavy,thick,-{Latex[length=2mm]}] (6.80,-0.02) -- (7.55,-0.02);
  % tầng gần trục
  \foreach \x/\lb in {0.78/BERT, 2.42/DeBERTa}{
    \fill[vnNavy] (\x,-0.02) circle (1.4pt);
    \node[below=1pt,text=vnNavy] at (\x,-0.02) {\lb};
  }
  % tầng xa trục
  \foreach \x/\lb in {1.56/RoBERTa, 3.88/{DeBERTaV3}}{
    \fill[vnNavy] (\x,-0.02) circle (1.4pt);
    \draw[vnNavy!45,line width=0.3pt] (\x,-0.02) -- (\x,-0.24);
    \node[below=0pt,text=vnNavy] at (\x,-0.24) {\lb};
  }
  % ---- đợt hiện đại hóa: hai mốc chỉ cách nhau 2 tháng nên nhãn phải tách ra,
  %      ModernBERT lùi hẳn sang trái, NeoBERT xuống tầng dưới cùng.
  \fill[vnNavy] (6.92,-0.02) circle (1.4pt);
  \draw[vnNavy!45,line width=0.3pt] (6.92,-0.02) -- (6.66,-0.16);
  \node[anchor=north east,text=vnNavy,align=right] at (6.64,-0.17) {ModernBERT\\12/2024};
  \fill[vnRed] (7.08,-0.02) circle (2.2pt);
  \draw[vnRed!45,line width=0.3pt] (7.08,-0.02) -- (7.08,-0.42);
  \node[below=0pt,text=vnRed,font=\tiny\bfseries,align=center] at (7.08,-0.42)
       {NeoBERT\\02/2025};
\end{tikzpicture}}
\end{frame}

% ---------------------------------------------------------------------
\begin{frame}{NeoBERT (2025) hiện đại hóa Encoder -- nhưng chỉ có tiếng Anh}
\vspace{0.3em}
\begin{columns}[T]
\begin{column}{0.52\textwidth}
  \textbf{NeoBERT} \cit{LeBreton2025} đưa các cải tiến của thế hệ LLM
  trở lại kiến trúc Encoder:
  \begin{itemize}\itemsep2pt
    \item \textbf{RoPE} -- mã hóa vị trí bằng phép quay
    \item \textbf{SwiGLU} -- mạng truyền thẳng có cổng
    \item \textbf{Pre-RMSNorm} -- chuẩn hóa trước mỗi lớp con
    \item \textbf{Sâu mà hẹp} -- 28 lớp, kích thước ẩn 768
  \end{itemize}

  \vspace{0.4em}
  Chỉ \textbf{250 triệu tham số} nhưng vượt BERT\textsubscript{large} và
  RoBERTa\textsubscript{large} trên chuẩn đánh giá MTEB, đồng thời hỗ trợ
  ngữ cảnh \textbf{4.096 token} -- gấp 8 lần BERT.
\end{column}

\begin{column}{0.45\textwidth}
  \small
  \begin{tabular}{@{}l l@{}}
    \toprule
    \rowcolor{vnLight}\multicolumn{2}{@{}l@{}}{\bfseries\color{vnNavy} Cấu hình NeoBERT}\\
    \midrule
    Số lớp / kích thước ẩn & 28 / 768\\
    Tham số                & $\approx$ 250 triệu\\
    Mã hóa vị trí          & RoPE\\
    Kích hoạt FFN          & SwiGLU\\
    Chuẩn hóa              & Pre-RMSNorm\\
    Ngữ cảnh tối đa        & 4.096 token\\
    Dữ liệu tiền huấn luyện& 2,1 nghìn tỷ token\\
    Ngôn ngữ               & \hl{chỉ tiếng Anh}\\
    \bottomrule
  \end{tabular}
  \vspace{4pt}
  {\footnotesize\color{vnCite} Bảng cấu hình theo \citet{LeBreton2025}}
\end{column}
\end{columns}
\end{frame}

% ---------------------------------------------------------------------
\begin{frame}{Bài toán: đưa NeoBERT sang tiếng Việt với chi phí chấp nhận được}
\vspace{0.2em}
\begin{columns}[T]
\begin{column}{0.49\textwidth}
  \textbf{Ba trở ngại}
  \begin{enumerate}\itemsep2pt
    \item Encoder tiếng Việt hiện có (PhoBERT \cit{Nguyen2020},
          ViDeBERTa \cit{Tran2023}) dùng kiến trúc cũ, giới hạn 512 token.
    \item Huấn luyện lại từ đầu cần \textbf{2,1 nghìn tỷ token} --
          bất khả thi với nhóm nghiên cứu nhỏ.
    \item Thay lớp véc-tơ rồi huấn luyện tiếp một cách ngây thơ dễ gây
          \term{mất thông tin}{catastrophic forgetting}.
  \end{enumerate}
\end{column}

\begin{column}{0.49\textwidth}
  \textbf{Ý tưởng của khóa luận}

  \vspace{2pt}
  \begin{block}{Tái sử dụng thay vì huấn luyện lại}
    Giữ nguyên \emph{phần thân} NeoBERT đã học từ tiếng Anh; chỉ thay
    lớp từ vựng, và khởi tạo lớp mới bằng cách \emph{mượn} không gian
    ngữ nghĩa tiếng Việt sẵn có của ViDeBERTa.
  \end{block}

  \vspace{2pt}
  \textbf{Mục tiêu cụ thể}
  \begin{itemize}\itemsep1pt
    \item Xây dựng \textbf{ViNeoBERT} bằng kỹ thuật SALT \cit{Lee2025}
          có cải tiến.
    \item Chứng minh bằng thực nghiệm rằng cách này \emph{tiết kiệm hơn
          nhiều} so với huấn luyện lại.
  \end{itemize}
\end{column}
\end{columns}
\end{frame}
